% GENERATED FILE. Edit canonical lessons, then run npm run generate:books.
% canonical-source-hash: fnv1a64:2065da4bb2aa527e
% canonical-lessons: AR-C212-khashin, AR-C212-lazij, AR-C212-hashsh, AR-C212-shaffaf, AR-C212-lami

\chapter{Rough, Sticky, Fragile, Transparent, Shiny}
\label{ch:ar-rough-sticky-fragile-transparent-shiny}
\hlchaptermodality{212}

\begin{chapteropening}
\textbf{By the end of this chapter:} \emph{I can say rough, sticky, fragile, transparent and shiny.}

Complete the last lesson of chapter 212: I can say rough, sticky, fragile, transparent and shiny.
\end{chapteropening}

\section[\texorpdfstring{\ar{خشن} (khashin)}{خشن (khashin)}]{\ar{خشن} (khashin) — rough}

\label{lesson:AR-C212-khashin}



\begin{quote}
Before the new one: say the Arabic for rare, then the Arabic for ordinary.
\end{quote}

\subsection*{You'll want to know: \ar{خشن}}
\textbf{\ar{خشن}} — \emph{khashin} — \textquotedblleft{}rough\textquotedblright{}.

\begin{grammarlens}[title={What you've built}]
One more word for this chapter.
\end{grammarlens}

\subsection*{Your turn}
\begin{itemize}
\raggedright
  \item \emph{Say it:} \emph{khashin}
  \item \emph{Say it:} \emph{khashin}, once more
  \item \emph{Say it:} say \emph{khashin}
  \item \emph{Recall:} say the Arabic for rare, then the Arabic for ordinary, then say \emph{khashin} again
\end{itemize}

\begin{tcolorbox}[breakable,colback=teal!4,colframe=teal!35!black,title={Before you move on}]
What does \ar{خشن} mean? (\textquotedblleft{}Rough\textquotedblright{}.) Say it once more.
\end{tcolorbox}
\section[\texorpdfstring{\ar{لزج} (lazij)}{لزج (lazij)}]{\ar{لزج} (lazij) — sticky}

\label{lesson:AR-C212-lazij}



\begin{quote}
Before the new one: say the Arabic for ordinary, then the Arabic for rough.
\end{quote}

\subsection*{You'll want to know: \ar{لزج}}
\textbf{\ar{لزج}} — \emph{lazij} — \textquotedblleft{}sticky\textquotedblright{}.

\begin{grammarlens}[title={What you've built}]
One more word for this chapter.
\end{grammarlens}

\subsection*{Your turn}
\begin{itemize}
\raggedright
  \item \emph{Say it:} \emph{lazij}
  \item \emph{Say it:} \emph{lazij}, once more
  \item \emph{Say it:} say \emph{lazij}
  \item \emph{Recall:} say the Arabic for ordinary, then the Arabic for rough, then say \emph{lazij} again
\end{itemize}

\begin{tcolorbox}[breakable,colback=teal!4,colframe=teal!35!black,title={Before you move on}]
What does \ar{لزج} mean? (\textquotedblleft{}Sticky\textquotedblright{}.) Say it once more.
\end{tcolorbox}
\section[\texorpdfstring{\ar{هش} (hashsh)}{هش (hashsh)}]{\ar{هش} (hashsh) — fragile}

\label{lesson:AR-C212-hashsh}



\begin{quote}
Before the new one: say the Arabic for rough, then the Arabic for sticky.
\end{quote}

\subsection*{You'll want to know: \ar{هش}}
\textbf{\ar{هش}} — \emph{hashsh} — \textquotedblleft{}fragile\textquotedblright{}.

\begin{grammarlens}[title={What you've built}]
One more word for this chapter.
\end{grammarlens}

\subsection*{Your turn}
\begin{itemize}
\raggedright
  \item \emph{Say it:} \emph{hashsh}
  \item \emph{Say it:} \emph{hashsh}, once more
  \item \emph{Say it:} say \emph{hashsh}
  \item \emph{Recall:} say the Arabic for rough, then the Arabic for sticky, then say \emph{hashsh} again
\end{itemize}

\begin{tcolorbox}[breakable,colback=teal!4,colframe=teal!35!black,title={Before you move on}]
What does \ar{هش} mean? (\textquotedblleft{}Fragile\textquotedblright{}.) Say it once more.
\end{tcolorbox}
\section[\texorpdfstring{\ar{شفاف} (shaffāf)}{شفاف (shaffāf)}]{\ar{شفاف} (shaffāf) — transparent}

\label{lesson:AR-C212-shaffaf}



\begin{quote}
Before the new one: say the Arabic for sticky, then the Arabic for fragile.
\end{quote}

\subsection*{You'll want to know: \ar{شفاف}}
\textbf{\ar{شفاف}} — \emph{shaffāf} — \textquotedblleft{}transparent\textquotedblright{}.

\begin{grammarlens}[title={What you've built}]
One more word for this chapter.
\end{grammarlens}

\subsection*{Your turn}
\begin{itemize}
\raggedright
  \item \emph{Say it:} \emph{shaffāf}
  \item \emph{Say it:} \emph{shaffāf}, once more
  \item \emph{Say it:} say \emph{shaffāf}
  \item \emph{Recall:} say the Arabic for sticky, then the Arabic for fragile, then say \emph{shaffāf} again
\end{itemize}

\begin{tcolorbox}[breakable,colback=teal!4,colframe=teal!35!black,title={Before you move on}]
What does \ar{شفاف} mean? (\textquotedblleft{}Transparent\textquotedblright{}.) Say it once more.
\end{tcolorbox}
\section[\texorpdfstring{\ar{لامع} (lāmiʿ)}{لامع (lāmiʿ)}]{\ar{لامع} (lāmiʿ) — shiny}

\label{lesson:AR-C212-lami}



\begin{quote}
Before the new one: say the Arabic for fragile, then the Arabic for transparent.
\end{quote}

\subsection*{You'll want to know: \ar{لامع}}
\textbf{\ar{لامع}} — \emph{lāmiʿ} — \textquotedblleft{}shiny\textquotedblright{}.

\begin{grammarlens}[title={What you've built}]
One more word for this chapter.
\end{grammarlens}

\subsection*{Your turn}
\begin{itemize}
\raggedright
  \item \emph{Say it:} \emph{lāmiʿ}
  \item \emph{Say it:} \emph{lāmiʿ}, once more
  \item \emph{Say it:} say \emph{lāmiʿ}
  \item \emph{Recall:} say the Arabic for fragile, then the Arabic for transparent, then say \emph{lāmiʿ} again
\end{itemize}

\begin{tcolorbox}[breakable,colback=teal!4,colframe=teal!35!black,title={Before you move on}]
What does \ar{لامع} mean? (\textquotedblleft{}Shiny\textquotedblright{}.) Say it once more.
\end{tcolorbox}
